\section{Type and Bilingual Support}
\label{sec:type}

One visual language, two font pipelines. English reports build with pdfLaTeX: T1
encoding plus \texttt{mathpazo}, so both text and mathematics are URW Palladio, a
Palatino clone. Chinese reports build with XeLaTeX: \texttt{ctex} handles line breaking
and punctuation, the text font becomes the OpenType TeX Gyre Pagella, and mathematics
stays in Pazo.\citep{ctex,texgyre}

The pipelines are not merged because Chinese has no second road. pdfLaTeX cannot break
CJK lines or load CJK fonts, and what comes out is a PDF with missing glyphs or arbitrary
breaks. The class therefore refuses \texttt{zh} under pdfLaTeX and tells you the command
to run. A hard error is much cheaper than a PDF that looks fine and has silently dropped
characters.

The two font names look like an inconsistency but are two implementations of one design:
Palladio is Type~1, Pagella is its OpenType redraw. The visible difference is nil; the
difference in coverage is not. Pagella carries →, ≥ and µ,
Palladio does not. Under the English pipeline, T1 leaves ≥ ≤ ≠ ≈ ∞ undefined and
typing them drops the characters \emph{silently}, so the class maps those five to math
symbols; arrows and unit signs come from \texttt{textcomp} and need no help.

\begin{table}[htbp]
  \centering
  \small
  \caption{The two font pipelines. Chinese also picks a character set: \texttt{windows} or
  \texttt{fandol}.}
  \label{tab:pipeline}
  \begin{tabular}{@{}lll@{}}
    \toprule
    Item & English & Chinese \\
    \midrule
    Engine     & pdfLaTeX           & XeLaTeX \\
    Text font  & URW Palladio (T1)  & TeX Gyre Pagella \\
    Math font  & Pazo Math          & Pazo Math \\
    CJK        & ---                & ctex, \texttt{scheme=plain} \\
    Encoding   & T1                 & TU + T1 \\
    Best for   & English, arXiv     & Chinese, mixed text \\
    \bottomrule
  \end{tabular}
\end{table}

Chinese also has to pick a character set. On Windows the system Song typeface looks
better, selected with \texttt{fontset=windows}; on any machine the bundled Fandol fonts
always compile, selected with \texttt{fontset=fandol}, at the cost of lighter strokes and
text that cannot be copied out of the PDF. CI takes the portable option, a personal
machine takes the pretty one. That is a stated trade-off, not an oversight.

Metadata lives in one place, the smallest convention in the template and the one that
saves the most time:

\begin{lstlisting}[language=TeX,
  caption={A typical \texttt{\textbackslash reportsetup}. Every key is optional.},
  label={lst:setup}]
\documentclass[logo]{boogonreport}
\reportsetup{
  title    = {A Technical Report Template: Design and Usage},
  author   = {BoogonClothman},
  number   = {TR-2026-01},
  version  = {v1.0.1},
  date     = {2026-10-22},
}
\end{lstlisting}

The same definition feeds the title block, the running head, the footer and the PDF
properties, so a renamed report cannot end up with stale metadata in the file.
